\let\negmedspace\undefined
\let\negthickspace\undefined
\documentclass[journal,onecolumn]{IEEEtran}

\usepackage[a5paper, margin=10mm]{geometry} % Safe with onecolumn
\usepackage{lmodern} % Better fonts (optional but recommended)
\usepackage{tfrupee} % Rupee symbol package

% Header fix
\setlength{\headheight}{12pt} % Safe value
\setlength{\headsep}{0pt}     % Reduce gap

\usepackage{gvv-book}
\usepackage{gvv}
\usepackage{cite}
\usepackage{amsmath,amssymb,amsfonts,amsthm}
\usepackage{algorithmic}
\usepackage{graphicx}
\usepackage{textcomp}
\usepackage{xcolor}
\usepackage{txfonts}
\usepackage{listings}
\usepackage{enumitem}
\usepackage{mathtools}
\usepackage{gensymb}
\usepackage{comment}
\usepackage[breaklinks=true]{hyperref}
\usepackage{tkz-euclide} 
\usepackage{listings}
% \usepackage{gvv}                                        
\def\inputGnumericTable{}                                 
\usepackage[latin1]{inputenc}                                
\usepackage{color}                                            
\usepackage{array}                                            
\usepackage{longtable}                                       
\usepackage{calc}                                             
\usepackage{multirow}                                         
\usepackage{hhline}                                           
\usepackage{ifthen}                                           
\usepackage{lscape}

\usepackage{amsmath,amssymb}
\usepackage{booktabs}
\usepackage{tikz}
\usetikzlibrary{arrows.meta,angles,quotes}

\begin{document}

\bibliographystyle{IEEEtran}
\vspace{3cm}

\title{4.10.16}
\author{EE25BTECH11014 - Bhoomika Lokesh}
% \maketitle
% \newpage
% \bigskip
{\let\newpage\relax\maketitle}

\renewcommand{\thefigure}{\theenumi}
\renewcommand{\thetable}{\theenumi}
\setlength{\intextsep}{10pt} % Space between text and floats

% Remove this line so equations use default numbering (1, 2, 3,...)
% \numberwithin{equation}{enumi}
\numberwithin{figure}{enumi}
\renewcommand{\thetable}{\theenumi}


\textbf{Question:} The area of the region bounded by the curve $y = x + 1$ 
and the lines $x=2$ and $x=3$ is:
\begin{multicols}{4}
\begin{enumerate}
    \item $\frac{7}{2} sq.units$
    \item $\frac{9}{2} sq.units$
    \item $\frac{11}{2} sq.units$
    \item $\frac{13}{2} sq.units$
\end{enumerate}
\end{multicols}
\textbf{Solution:}
y = x + 1 is
\begin{align}
 \vec{n^Tx}=c\\
\vec{n^T}=\myvec{1\\-1}
\end{align}

where $\vec{n^T}$ is the normal vector of the line.\\
Let the points of intersection of the curve $y = x + 1$ and the lines $x=2$ and $x=3$ be
$\vec{A},\vec{B}.$\\
For $\vec{A}$:
\begin{align}
\myvec{1 & 0 \\ 
-1 & 1}
\myvec{x\\y}=\myvec{2\\1}
\end{align}
The augmented matrix of the above is:
\begin{align}
\myvec{\begin{array}{cc|c}
1 & 0 & 2 \\ 
-1 & 1 & 1
\end{array}}.
\end{align}
Then, 
\begin{align}
    \vec{R_1}\longrightarrow\vec{R_2}+\vec{R_1}\\
  \myvec{\begin{array}{cc|c}
1 & 0 & 2 \\ 
0 & 1 & 3
\end{array}}\\
\therefore \myvec{x\\y}=\myvec{2\\3}.
\end{align}

For $\vec{B}$:
\begin{align}
\myvec{1 & 0 \\ 
-1 & 1}
\myvec{x\\y}=\myvec{3\\1}
\end{align}

The augmented matrix of the above is:
\begin{align}
\myvec{\begin{array}{cc|c}
1 & 0 & 3 \\ 
-1 & 1 & 1
\end{array}}.
\end{align}

Then, 
\begin{align}
    R_2 \longrightarrow R_2 + R_1 \\
  \myvec{\begin{array}{cc|c}
1 & 0 & 3 \\ 
0 & 1 & 4
\end{array}}\\
\therefore  \myvec{x\\y}=\myvec{3\\4}.
\end{align}
Hence the points of the quadrilateral are $\brak{2,0},\brak{2,3},\brak{3,4}$ and $\brak{3,0}$ (which are $\vec{A},\vec{B},\vec{C}$ and $\vec{D}$ respectively.)

 Area of the quadrilateral can be calculated by using the formula,
 \begin{align}
 Ar=\frac{1}{2}\left\|\brak{\vec{A}-\vec{C}}\times\brak{\vec{D}-\vec{B}}\right\|\\
 \brak{\vec{A}-\vec{C}}=\myvec{-1\\-4}   \\\brak{\vec{B}-\vec{D}}=\myvec{1\\3}  \\
 \brak{\vec{A}-\vec{C}}\times\brak{\vec{D}-\vec{B}}=\begin{vmatrix}-1&1 \\-4&3\end{vmatrix}=-3-4=-7\\
 \therefore Ar=\frac{1}{2}\left|-7\right|=\frac{7}{2}.
 \end{align}
 Therefore (option 1) $\frac{7}{2} sq.units$ is the correct answer.

\begin{figure}[H]
\begin{center}
\includegraphics[width=0.7\columnwidth]{figs/fig_8.png}
\end{center}
\caption{Graph}
\label{fig:fig.py}
\end{figure}

 
\end{document}